%%%%%%%%%%%%%%%%%%%%%%%%%%%%%%%%%%%%%%%%%%%%%%%%%%%%%%%%%%%%%%%%%%%%%%%%%%%
%%  conclusion.tex  –  Заключение (ненумерованная глава)
%%  v9 – Последовательное изложение всех защищаемых результатов
%%       по цепочке: НВ динамика → распределённость → эффективность →
%%       адаптация → сертификация → предметная область → платформа.
%%       Соответствует структуре из 5 глав.
%%%%%%%%%%%%%%%%%%%%%%%%%%%%%%%%%%%%%%%%%%%%%%%%%%%%%%%%%%%%%%%%%%%%%%%%%%%

\chapter*{ЗАКЛЮЧЕНИЕ}
\addcontentsline{toc}{chapter}{Заключение}
\markboth{ЗАКЛЮЧЕНИЕ}{ЗАКЛЮЧЕНИЕ}

\section*{Краткая характеристика основных результатов
(восемь утверждений по восьми задачам исследования)}

Основные научные результаты диссертационного исследования
систематизированы ниже в виде восьми утверждений в строгом
соответствии восьми поставленным задачам исследования
(см.~Введение). Каждое утверждение фиксирует, как поставленная
задача решена в работе и где её решение представлено.

\begin{enumerate}[leftmargin=2em, itemsep=0.3em]

\item \textit{По задаче <<Систематический обзор подходов,
применяемых в современных гибридных системах обнаружения и
предотвращения вторжений, и выявление их недостатков и
уязвимостей>>:} в Главе~\ref{ch:literature_review} проведён
систематический обзор подходов, применяемых в современных
Hybrid~IDPS; задокументированы критические отказы моделей в
здравоохранении (падение AUC до 25\,\%), финансах
(25{,}7\,\%~$\to$~68{,}6\,\% ошибок), промышленности
(94{,}2\,\%~$\to$~30\,\% точности) и кибербезопасности
($-$59\,\% обнаружения), а также отставание устойчивости
моделей на 25--65~п.\,п.\ в условиях состязательных воздействий.

\item \textit{По задаче <<Обзор методов повышения робастности к
состязательным атакам и обеспечения вычислительной
эффективности>>:} систематизированы пять семейств моделей и
алгоритмов Hybrid~IDPS (статистические, классические ML,
глубокое обучение, гибридные, теоретико-игровые;
см.~рис.~\ref{fig:orga_models_idps}) и проведён сравнительный
анализ ИИ-усиленных коммерческих продуктов (CrowdStrike Charlotte
AI, SentinelOne Purple AI, Microsoft Security Copilot, Google
Mandiant+Gemini, Cisco Hypershield, IBM QRadar AI, Palo~Alto
Cortex~XSIAM), что зафиксировало отсутствие единой архитектуры,
обеспечивающей одновременно сертифицированную устойчивость,
распределённую приватность, вычислительную эффективность и
калиброванную точность на динамических графовых структурах.

\item \textit{По задаче <<Формализовать математическую модель
гибридной системы обнаружения и предотвращения вторжений и
поставить задачу повышения её робастности и эффективности>>:}
в Главе~\ref{ch:methods} сформулирована единая кортежная
математическая модель Hybrid~IDPS, дано формальное определение
шести свойств П1--П6 и матрицы <<условие\,---\,требование>> из
12~ячеек, и в явном виде поставлена многокритериальная задача
повышения робастности, точности, эффективности и приватности.

\item \textit{По задаче <<Разработать алгоритмы и модели
повышения робастности и эффективности гибридной системы
обнаружения вторжений в распределённых и нестационарных
условиях>>:} разработаны семь моделей М1--М7, реализующие
непрерывно-временную графовую динамику с сертификатом по
Липшицу/Грёнуоллу (\textbf{CT-TGNN/SDE-TGNN}; 99{,}0\,\% F1
на CIC-IoT-2023, сертифицированный радиус $\rho=0{,}25$),
византийско-устойчивое федеративное обучение с дифференциальной
приватностью (\textbf{FedLLM-API}; 87{,}1\,\% F1 при 40\,\%
повреждённых клиентов, $\varepsilon=0{,}85$), многоуровневые
темпоральные представления (\textbf{TripleE-TGNN}; 96{,}8\,\% F1
на многостадийных атаках), модель пространства состояний с
устойчивостью к отравлению (\textbf{MambaShield};
$\mathcal{O}(L\log L)$, 91{,}4\,\% F1 при 25\,\% отравлении,
0{,}5~мс латентности), стохастический трансформер с UC-HGP
(ECE\,=\,0{,}078; $-43\,\%$ нагрузки SOC) и единую
теоретико-игровую систему сертификации
(\textbf{M7-Certification}; 94{,}7\,\% F1 против адаптивных
противников).

\item \textit{По задаче <<Разработать программное обеспечение,
реализующее предложенные модели и алгоритмы (формирование
требований, сценариев использования, архитектуры и программной
реализации)>>:} в Главе~\ref{ch:platform} разработана
программная платформа \textbf{RobustIDPS.ai}
(DOI~10.5281/zenodo.19129512), реализующая модели М1--М7 в
составе единой системы: 17~нейросетевых моделей,
10~функциональных групп, 51~страница прогрессивного
веб-приложения (PWA), четырёхуровневая система защиты больших
языковых моделей CyberSecLLM (94{,}5\,\% обнаружения по
517~сценариям атак, 89{,}1\,\% zero-shot F1 на наборе
кибербезопасности, +20{,}5~п.\,п.\ над GPT-4); сформированы
функциональные требования, сценарии использования и архитектура
системы (FastAPI/PyTorch~2.x/PyTorch~Geometric, React~18/PWA,
PostgreSQL~16/Redis~7, Docker~Compose; четыре LLM-провайдера и
плагин Suricata EVE-JSON; см.~\S\ref{sec:tech_stack_ui}).

\item \textit{По задаче <<Провести экспериментальные исследования
разработанных алгоритмов и моделей на наборах данных-эталонах
(benchmarks)>>:} в Главе~\ref{ch:experiments} выполнена
комплексная экспериментальная программа на шести наборах данных
сетевого трафика (CIC-IoT-2023, CSE-CIC-IDS2018, UNSW-NB15,
MS-OUEDA, Container Security, Edge-IoT), включающая 84{,}2~млн
помеченных образцов, 84 типа атак и 23~метрики качества,
робастности, эффективности, приватности и калибровки;
подтверждены все 12~ячеек матрицы <<условие--требование>> с
превосходством над базовыми методами на 2{,}9--5{,}6~п.\,п.

\item \textit{По задаче <<Интегрировать разработанное программное
обеспечение в реальный конвейер анализа сетевого трафика>>:}
получены \textbf{два акта внедрения}, подтверждающие интеграцию
разработанного программного комплекса RobustIDPS.ai в реальные
производственные конвейеры анализа сетевого трафика и защиты
сетевой инфраструктуры: от ООО~<<НИИ~ИКТ>> (г.~Новосибирск,
от 03.06.2026~г.; интеграция моделей, устойчивых к атакам
отравления, в действующую систему обнаружения вторжений и
активное использование платформы RobustIDPS.ai штатными
специалистами подразделений кибербезопасности) и от
ООО~<<Ксайрикс>> (г.~Москва, 2026~г.; интеграция в существующее
программное обеспечение системы анализа сетевого трафика и
защиты сетевой инфраструктуры~--- обнаружение многостадийных
атак, выявление ранее неизвестных паттернов и адаптация к
дрейфу).

\item \textit{По задаче <<Продемонстрировать, что разработанное
программное обеспечение допускает воспроизведение в другой
практической области (на примере атак на БПЛА)>>:} в
Главе~\ref{ch:uav} реализована трёхъярусная архитектура
<<борт\,---\,дронопорт\,---\,облако>> с применением М1--М7 для
защиты бортовых систем компьютерного зрения беспилотных
летательных аппаратов: бортовой ярус выполняет инференс
MambaShield с латентностью 0{,}5~мс/поток, дронопорт агрегирует
обновления через FedLLM-API, облачный ярус реализует
теоретико-игровую сертификацию M7. Подтверждена доменная
независимость разработанного математического и алгоритмического
каркаса; третий акт внедрения (Центр искусственного интеллекта
НИЯУ МИФИ, г.~Москва) находится в стадии оформления и закрепит
эту воспроизводимость документально.

\end{enumerate}

\noindent Развёрнутое содержание перечисленных результатов
приводится в последующих разделах настоящего Заключения.

\bigskip

В настоящей диссертации решена фундаментальная задача обеспечения
робастности, эффективности, точности и приватности \emph{гибридных
систем обнаружения и предотвращения вторжений} (Hybrid~IDPS),
функционирующих в состязательных, распределённых и нестационарных
условиях. Исследование было мотивировано задокументированными
отказами существующих IDPS-решений в здравоохранении, финансах,
промышленном IoT, автономных системах и кибербезопасности, где
робастность систем отставала от точности обнаружения на
25--65~процентных пунктов. Единая математическая модель
Hybrid~IDPS была формализована в виде семёрки
$\mathcal{H} = (\mathcal{V}_t, \mathcal{E}_t, X_t, A_t, f_\theta,
g_\psi, \pi)$, и шесть сквозных свойств системы (состязательная
устойчивость, робастность, эффективность, совместимость с
распределённой средой, человеческий фактор, нестационарность)
систематически декомпозированы на семь подзадач, решённых семью
новыми моделями и алгоритмами~М1--М7. Графовые нейронные сети
использованы в работе как одна из ветвей семейства моделей,
применяемых в Hybrid~IDPS, наряду со статистическими, классическими
ML-, гибридными и теоретико-игровыми подходами. В соответствии с
темой диссертации вклад работы позиционирован как \emph{разработка
состязательных устойчивых моделей на основе искусственного
интеллекта в гибридных системах обнаружения и предотвращения
вторжений для сетевой безопасности}: сектор защиты сетевого
трафика (NIDPS) выступает основной инстанцией развёртывания, а
сектор защиты бортовых систем БПЛА~--- подтверждением переноса
каркаса AI-моделей в смежную предметную область. Получены
следующие результаты.


%%---------------------------------------------------------------------
\section*{Основные результаты}
%%---------------------------------------------------------------------

\noindent\textbf{1. CT-TGNN и SDE-TGNN: Сертифицированная динамика графов непрерывного времени (Пробел~1).}
Разработана первая темпоральная графовая нейронная сеть непрерывного времени, объединяющая графовую передачу сообщений с динамикой нейронных ОДУ на эволюционирующих топологиях. Обучение методом сопряжённых чувствительностей достигает $O(1)$ памяти; динамика с ограниченной константой Липшица обеспечивает сертифицированную робастность через неравенство Грёнуолла. Стохастическое расширение (SDE-TGNN) заменяет детерминированную динамику стохастическими дифференциальными уравнениями Ито и выводит аналитическое распространение моментов из уравнения Фоккера--Планка. Экспериментальная валидация на трёх бенчмарках общим объёмом 52 миллиона образцов демонстрирует средневзвешенный F1 99,0\%, ожидаемую ошибку калибровки 0,042, сертифицированную состязательную робастность 76,2\% при радиусе возмущения $R=0.25$ и кросс-доменный перенос 90,5\% F1 без дообучения.

\noindent\textbf{2. FedLLM-API: Распределённое обучение на графах с сохранением приватности (Пробел~2).}
Разработан оригинальный византийско-устойчивый федеративный оптимизатор для графов, сочетающий покоординатную агрегацию усечённым средним с $(\varepsilon{=}0.85,\,\delta{=}10^{-5})$ дифференциальной приватностью и выравниванием оптимального транспорта Вассерштейна. Сходимость со скоростью $\mathcal{O}(1/\sqrt{T})$ доказана при 30\% повреждённых участников. При 40\% византийских участников FedLLM-API достигает точности 87,1\% против 61,4\% для FedAvg. Интеграция LLM обеспечивает 87,6\% F1 при обнаружении новых категорий атак в режиме zero-shot.

\noindent\textbf{3. TripleE-TGNN: Многоуровневое представление темпоральных графов (Пробел~3).}
Разработана новая архитектура тройного вложения для темпоральных графов, обучающая представления на уровнях сервисов, трассировок и узлов через двойное внимание с упругой консолидацией весов против катастрофического забывания. Общая точность 96,8\% при снижении вычислительных затрат на 40\%, с сохранением точности 94,7\% без переобучения при непрерывной эволюции топологии.

\noindent\textbf{4. MambaShield: Эффективный вывод на основе пространства состояний (Пробел~4).}
Разработана новая селективная архитектура пространства состояний, снижающая сложность вывода с $\mathcal{O}(L^{2})$ до $\mathcal{O}(L\log L)$ через свёртки БПФ с прогрессивной состязательной дистилляцией. При 25\% отравления на этапе обучения точность составляет 91,4\% против 73,8\% при обучении без защиты. PAC-байесовские сертификаты обобщения подтверждают границу с запасом 2,7\%. Задержка вывода 0,5\,мс на поток является наименьшей среди всех моделей.

\noindent\textbf{5. Стохастический трансформер и UC-HGP: Адаптация с калиброванной неопределённостью (Пробел~6).}
Разработан первый вывод ГНС с калиброванной неопределённостью, сочетающий вариационное байесовское внимание с иерархическими гауссовскими процессами, достигающий декомпозиции эпистемической и алеаторной неопределённости с ECE 0,078 (снижение на 59\% по сравнению с детерминированным вниманием) и 43\% снижения нагрузки аналитиков через сортировку на основе неопределённости.

\noindent\textbf{6. Теоретико-игровая сертификация: Единая система робастности (Пробел~7).}
Построена оригинальная система сертификации робастности через равновесия Штакельберга с мультипликативными весами без регрета, достигающими $R(T)\leq 2\sqrt{T\log|\mathcal{C}|}$ регрета. Регуляризация согласованности связывает все семь методов с доказанной сходимостью. Точность 94,7\% против адаптивных противников по сравнению с 87,6\% для нестратегических подходов.

\noindent\textbf{7. CyberSecLLM: Предметно-ориентированная фундаментальная модель (Пробел~5).}
Разработана новая предметно-ориентированная фундаментальная модель с блоками селективного пространства состояний, дополненными ОДУ, достигающая средней точности 89,1\% в режиме zero-shot на бенчмарке из 11 задач CyberSecBench --- на 20,5 п.\,п.\ выше GPT-4 --- при обработке графовых последовательностей событий длиной в миллион токенов за 2,3 секунды.

\noindent\textbf{8. Предметная адаптация.}
Все методы реализованы в области кибербезопасности через формальный конвейер преобразования данных в граф, унифицированную таксономию атак из 34 классов, конвейер инженерии 83 признаков и предметно-ориентированные конфигурации гиперпараметров. Адаптация демонстрирует переносимость математической основы в конкретную среду развёртывания.

\noindent\textbf{9. Экспериментальная валидация.}
Комплексные эксперименты на шести бенчмарк-наборах данных общим объёмом 84,2 миллиона размеченных образцов, 84 категориях атак и 23 метриках оценки подтверждают, что единая система превосходит лучшие базовые линии на 2,9--5,6 процентных пунктов. Абляционные исследования подтверждают, что каждый компонент вносит измеримый вклад. Состязательная робастность продемонстрирована при шести семействах атак. Кросс-доменный перенос на речь и здравоохранение подтверждает предметную независимость.

\noindent\textbf{10. Платформа RobustIDPS.ai.}
Полная система реализована как платформа с открытым исходным кодом (DOI: 10.5281/zenodo.19129512), содержащая 51 страницу, 10 функциональных групп и 17 интегрированных моделей. Четырёхуровневая система защиты LLM достигает обнаружения 94,5\% по 517 сценариям атак с четырьмя коммерческими бэкендами LLM. Модуль Multi-Agent PQC-IDS расширяет обнаружение на постквантовые протокольные среды. Валидация пользовательским исследованием подтверждает 34\% снижение времени сортировки аналитиками. Платформа функционирует как плагин для Snort, Suricata и аналогичных систем СОПСВ.


%%---------------------------------------------------------------------
\section*{Подтверждение матрицы условие--требование}
%%---------------------------------------------------------------------

Экспериментальная программа подтверждает, что все двенадцать пересечений матрицы условие--требование $3\times4$ покрыты с подтверждённой производительностью:

\begin{itemize}
\item Состязательность $\times$ Робастность: сертификаты Липшица CT-TGNN (98,3\% точности)
\item Состязательность $\times$ Точность: калибровка стохастического трансформера (ECE 0,078)
\item Состязательность $\times$ Эффективность: вывод MambaShield $\mathcal{O}(L\log L)$ (0,5\,мс/поток)
\item Состязательность $\times$ Приватность: теоретико-игровая сертификация с композицией ДП
\item Нестационарность $\times$ Робастность: сертифицированная робастность SDE-TGNN (76,2\%)
\item Нестационарность $\times$ Точность: многоуровневые вложения TripleE-TGNN (96,8\%)
\item Нестационарность $\times$ Эффективность: прогрессивная дистилляция MambaShield (91,4\% при отравлении)
\item Нестационарность $\times$ Приватность: прямо совместимое обнаружение PQC (96,2\%)
\item Распределённость $\times$ Робастность: византийская устойчивость FedLLM-API (87,1\% при 40\% повреждении)
\item Распределённость $\times$ Точность: FedLLM-API + LLM zero-shot (87,6\% F1)
\item Распределённость $\times$ Эффективность: ускоренная сходимость ОТ (на 35\% быстрее)
\item Распределённость $\times$ Приватность: $(\varepsilon{=}0.85,\delta{=}10^{-5})$ дифференциальная приватность
\end{itemize}


%%---------------------------------------------------------------------
\section*{Промышленная валидация}
%%---------------------------------------------------------------------

Получены два акта внедрения, подтверждающие практическую применимость разработанных состязательных устойчивых моделей на основе искусственного интеллекта для Hybrid~IDPS в задачах сетевой безопасности: ООО <<Научно-исследовательский институт информационно-коммуникационных технологий>> (ООО <<НИИ ИКТ>>, г.~Новосибирск; интеграция моделей обнаружения аномалий сетевого трафика, устойчивых к атакам отравления, в действующую систему обнаружения вторжений организации и активное использование программного комплекса RobustIDPS.ai штатными специалистами подразделений кибербезопасности) и ООО <<Ксайрикс>> (г.~Москва; применение разработанных решений в составе системы анализа сетевого трафика и защиты сетевой инфраструктуры~--- обнаружение многостадийных атак и координированных аномалий в графе сетевого взаимодействия, выявление ранее неизвестных паттернов аномального поведения на основе анализа неопределённости и адаптация к изменению статистических характеристик трафика). Акт от Центра искусственного интеллекта Национального исследовательского ядерного университета <<МИФИ>>, г.~Москва, подтверждающий применение разработанного каркаса AI-моделей Hybrid~IDPS к защите бортовых систем БПЛА (Глава~\ref{ch:uav}), находится в стадии оформления.


%%---------------------------------------------------------------------
\section*{Направления дальнейших исследований}
%%---------------------------------------------------------------------

Результаты настоящей диссертации открывают ряд направлений для дальнейших исследований.

Во-первых, расширение динамики графов непрерывного времени на гетерогенные темпоральные гиперграфы позволит моделировать взаимодействия высшего порядка (например, многосторонние транзакции, групповые коммуникации), которые невозможно представить попарными рёбрами, с сохранением системы сертификации по Липшицу.

Во-вторых, интеграция разработанных методов с большими мультимодальными фундаментальными моделями (зрение--язык--граф) позволит совместно обрабатывать графы сетевого трафика с трассировками системных вызовов, телеметрией конечных точек и потоками DNS-запросов для целостного обнаружения угроз.

В-третьих, развёртывание квантованных моделей Multi-Agent PQC-IDS на микроконтроллерах ARM Cortex-M и RISC-V расширит классификацию постквантового трафика на ресурсно-ограниченные периферийные устройства IoT.

В-четвёртых, применение методов формальной верификации (на основе SMT или абстрактной интерпретации) для доказательства ограниченных свойств безопасности конвейера защиты LLM усилит гарантии сертификации за пределами эмпирической оценки.

В-пятых, расширение системы федеративного обучения на полностью децентрализованные (одноранговые) условия без доверенного сервера агрегации устранит единую точку отказа в текущей звёздной топологии.

В-шестых, адаптация полной системы к дополнительным областям --- точное земледелие (пространственно-временные графы датчиков посевов), автономное вождение (динамические графы слияния датчиков) и надзор за финансовыми рынками (темпоральные гиперграфы транзакций) --- подтвердит предметную независимость, заявленную в теоретическом анализе.

Эти направления естественно развивают математические основы, алгоритмическую инфраструктуру и программную платформу, разработанные в настоящей диссертации, и в совокупности направлены на продвижение уровня развития робастного и надёжного ИИ для динамических графовых данных.
